%=====================================================================
%  Advanced Optimization (Fall 2026) -- Midterm Examination (exam only; model answers kept out of this public repo)
%  Nocedal & Wright, Numerical Optimization, 2nd ed., Chapters 2-6
%  Five questions, one per chapter, no sub-parts. Each question is anchored in a
%  small example that is not in the book, so the book helps but does not answer it.
%=====================================================================
\documentclass[11pt,a4paper]{article}
\usepackage[margin=2.3cm]{geometry}
\usepackage{amsmath,amssymb}
\usepackage{xcolor}
\usepackage{enumitem}
\setlist[enumerate,1]{label=\textbf{A\arabic*.}, leftmargin=2.6em, itemsep=2pt}
\usepackage{fancyhdr}

\newcommand{\R}{\mathbb{R}}
\newcommand{\T}{^{\mathsf T}}
\newcommand{\gradf}{\nabla f}
\newcommand{\hessf}{\nabla^{2} f}
\DeclareMathOperator{\diag}{diag}

\newcommand{\versionlabel}{}

\pagestyle{fancy}
\fancyhf{}
\lhead{Advanced Optimization --- Midterm \versionlabel}
\rhead{Fall 2026, Sogang University}
\cfoot{\thepage}

\begin{document}

\begin{center}
  {\LARGE\bfseries Advanced Optimization --- Midterm Examination \versionlabel}\\[4pt]
  Fall 2026 \quad $\cdot$ \quad 75 minutes \quad $\cdot$ \quad Part A: 20 true/false \quad $\cdot$ \quad Part B: 4 questions
\end{center}

\medskip
\noindent No calculator is needed; leave fractions and square roots as they are.
\textbf{Part A} (40 points): twenty statements, 2 points each. Mark each T or F. If F, give a counterexample or a one-line reason; a bare F earns 1 point, a bare T earns 2 if correct.
Do Part A first.

\noindent\textbf{Part B} (60 points, 15 each): Questions 1 and 2 ask for explanations; Questions 3 and 4 ask you to work a small example and then to explain what it shows.
The credit is for the explanation in your own words; a quoted passage without explanation earns none.

\bigskip
\section*{Part A --- True or false}

\textbf{Chapter 2}
\begin{enumerate}
\item If $x^{*}$ is a strict local minimizer of $f$, then it is an isolated local minimizer.

\item If $\gradf(x^{*})=0$ and $\hessf(x^{*})$ is positive semidefinite, then $x^{*}$ is a local minimizer.

\item Newton's method is invariant under a linear change of variables $x=Sz$ ($S$ nonsingular): the iterates in $z$ map to the iterates in $x$. Steepest descent is not.

\item The sequence $x_{k}=1+k^{-k}$ converges to $1$ Q-superlinearly but not Q-quadratically.
\end{enumerate}

\textbf{Chapter 3}
\begin{enumerate}[resume]
\item Let $f$ be smooth and bounded below along the ray $\{x_{k}+\alpha p_{k}:\alpha>0\}$, with $p_{k}$ a descent direction. Then for any $0<c_{1}<c_{2}<1$ there exist step lengths satisfying the strong Wolfe conditions.

\item Zoutendijk's theorem implies that steepest descent with a Wolfe line search converges to a local minimizer.

\item Let $x_{k}\to x^{*}$ with $\gradf(x^{*})=0$, $\hessf(x^{*})\succ0$, and let $p_{k}$ be the Newton direction. With a Wolfe line search with $c_{1}\le\tfrac12$, the unit step $\alpha_{k}=1$ satisfies the Wolfe conditions for all sufficiently large $k$.

\item For $f(x)=\tfrac12x\T Qx$ with $Q\succ0$, steepest descent with exact line searches reaches the minimizer in at most $n$ iterations.
\end{enumerate}

\textbf{Chapter 4}
\begin{enumerate}[resume]
\item In Algorithm 4.1, the trust-region radius is increased whenever $\rho_{k}>\tfrac34$.

\item The Cauchy point always lies on the boundary of the trust region.

\item If $B_{k}$ is indefinite, the exact solution of the trust-region subproblem lies on the boundary $\|p\|=\Delta_{k}$.

\item The dogleg method for the trust-region subproblem requires $B_{k}$ to be positive definite.
\end{enumerate}

\textbf{Chapter 5}
\begin{enumerate}[resume]
\item Applied to $Ax=b$ with $A\in\R^{n\times n}$ symmetric positive definite, the conjugate gradient method needs exactly $n$ iterations to reach the solution in exact arithmetic.

\item In the conjugate gradient method the residuals $r_{0},r_{1},\dots$ are mutually orthogonal and the directions $p_{0},p_{1},\dots$ are mutually conjugate.

\item The $k$-th CG iterate $x_{k}$ minimizes $\phi(x)=\tfrac12x\T Ax-b\T x$ over the whole affine set $x_{0}+\mathrm{span}\{p_{0},\dots,p_{k-1}\}$, not only along the last direction $p_{k-1}$.

\item The number of CG iterations needed to reach a given accuracy is determined by the condition number $\kappa(A)$.
\end{enumerate}

\textbf{Chapter 6}
\begin{enumerate}[resume]
\item If $f$ is strongly convex, then $s_{k}\T y_{k}>0$ for any two points $x_{k},x_{k+1}$, whatever line search is used.

\item If $B_{k}$ is positive definite and $s_{k}\T y_{k}>0$, then the BFGS update $B_{k+1}$ is positive definite.

\item If $B_{k}$ is positive definite and $s_{k}\T y_{k}>0$, then the SR1 update $B_{k+1}$ is positive definite.

\item BFGS with a Wolfe line search satisfies $\liminf\|\gradf_{k}\|=0$ for every twice continuously differentiable $f$ that is bounded below.
\end{enumerate}

\section*{Part B}

\section*{Question 1 (Chapter 3)}

What is the Armijo (sufficient decrease) condition? What are the Wolfe conditions?
Why is the Armijo condition with backtracking not always enough --- when is it enough, and when is it not?


\section*{Question 2 (Chapters 3--6)}

In unconstrained minimization the objective is modeled locally by a quadratic $m_{k}(p)=f_{k}+\gradf_{k}\T p+\tfrac12p\T B_{k}p$. Positive definiteness of $B_{k}$ (or of the Hessian) is often assumed, but not always needed.
\begin{itemize}
\item[(A)] List the methods of Chapters 3--6 that require $B_{k}$ to be positive definite, and say for each what breaks if it is not.
\item[(B)] List the methods that do not require it, and explain in each case how that is possible.
\end{itemize}


\section*{Question 3 (Chapter 4)}

Consider the model $m(p)=g\T p+\tfrac12p\T Bp$ with $g=(1,0)\T$ and $B=\diag(1,-1)$, and the trust region $\|p\|\le1$.
A line search method would first make $B$ positive definite (e.g.\ $B+\tau I$ with $\tau>1$) and then search along $-(B+\tau I)^{-1}g$: find the best model value such a method can reach, for any $\tau$ and any step length within the trust region.
Then solve the trust-region subproblem exactly using Theorem 4.1 (note what happens to the condition $\|p(\lambda)\|=\Delta$), give its model value, and explain what the trust-region step does with the negative curvature that the line search step cannot do.


\section*{Question 4 (Chapter 5)}

Let $A=\diag(1,3)$, $b=0$, $x_{0}=(3,1)\T$, with residual $r_{k}=Ax_{k}-b$; the solution is $x^{*}=0$.
Starting from $p_{0}=-r_{0}$, take the exact line search step to $x_{1}$ and compute $r_{1}$.
Then find the direction of the form $p_{1}=-r_{1}+\beta p_{0}$ that is conjugate to $p_{0}$ (i.e.\ $p_{0}\T Ap_{1}=0$), take the exact line search step along it, and verify that you arrive at $x^{*}$.
For comparison, compute where exact-line-search steepest descent ($p_{1}=-r_{1}$) lands instead, and explain why it misses $x^{*}$ while the conjugate direction does not.


\end{document}
